%\documentclass[12pt,oneside]{amsart}
\documentclass{scrartcl}

%\documentclass{article}

\usepackage{amsthm,amsmath,amssymb}
\usepackage[numbers]{natbib}

%\usepackage[colorlinks]{hyperref}
\usepackage{hyperref}
\hypersetup{
    colorlinks,%
    citecolor=black,%
    filecolor=black,%
    linkcolor=black,%
    urlcolor=black
}
\usepackage{hypernat}
\usepackage[top=2.5cm, bottom=2.5cm, left=2.8cm,
right=2.8cm]{geometry}
\usepackage{graphicx}


%\setlength{\parskip}{0pt}
%\setlength{\parsep}{0pt}
%\setlength{\headsep}{0pt}
%\setlength{\topskip}{0pt}
%\setlength{\topmargin}{0pt}
%\setlength{\topsep}{0pt}
%\setlength{\partopsep}{0pt}

\linespread{1}

%\usepackage{marvosym}

%\textwidth = 6.5 in \textheight = 9.2 in \oddsidemargin = 0.0 in
%\evensidemargin = 0.0 in \topmargin = -0.0 in \headheight = 0.0 in
%\headsep = 0.0in \parskip = 0.0in \parindent = 0.0in
%\def\baselinestretch{0.871}


\newtheorem{theorem}{Theorem}[section]
\newtheorem{proposition}[theorem]{Proposition}
\newtheorem{problem}[theorem]{Problem}
\newtheorem{fact}[theorem]{Fact}
\newtheorem{lemma}[theorem]{Lemma}
\newtheorem{defn}[theorem]{Definition}
\newtheorem{corollary}[theorem]{Corollary}
\newtheorem{remark}{Remark}[section]
\newtheorem{example}[theorem]{Example}

\newcommand{\E}{{\mathbb E}}
\newcommand{\se}{{\mathcal{E}}}
\newcommand{\R}{{\mathbb R}}
\renewcommand{\P}{{\mathbb P}}
\newcommand{\ac}{{\mathcal{A}}}
\newcommand{\A}{{\mathcal{A}}}
\newcommand{\B}{{\mathcal{B}}}
\newcommand{\C}{{\mathcal{C}}}
\newcommand{\M}{{\mathcal{M}}}
\newcommand{\Mf}{{\mathfrak{M}}}
\newcommand{\T}{{\mathcal{T}}}
\newcommand{\Z}{{\mathcal{Z}}}
\newcommand{\K}{{\mathcal{K}}}
\newcommand{\lc}{{\mathcal{L}}}
\newcommand{\Ps}{{\mathcal{P}}}
\newcommand{\G}{{\mathcal{G}}}
\newcommand{\pa}{{\mathring{p}}}
\newcommand{\F}{{\cal F}}
\newcommand{\samp}{{\mathcal{X}}}
\newcommand{\X}{{\mathcal{X}}}
\newcommand{\hi}{{\hat{\Phi}}}
\newcommand{\data}{{\text{data}}}

\newcounter{rcnt}[section]
\renewcommand{\thercnt}{(\roman{rcnt})}

\def\qt#1{\qquad\text{#1}}

\def\argmin{\mathop{\rm argmin}}
\def\argmax{\mathop{\rm argmax}}


\setlength{\parskip}{1.4 \medskipamount}
%\sloppy
%\linespread{1.3}

\begin{document}

\title{STAT 238 - Bayesian Statistics  \\
  Lecture Ten} 
\subtitle{Spring 2026, UC Berkeley}

\author{Aditya Guntuboyina}


\date{11 Feb 2026}

\maketitle

\section{Additional Comments on the Kidney Cancer Data Analysis}
For the kidney cancer dataset $(X_i, n_i), i = 1, \dots, N$ ($N$
denotes the total number of counties, $n_i$ is the average population
of county $i$ and $X_i$ is the number of deaths due to kidney cancer
in the fixed time period $1980-89$). Consider the following three
models for this dataset:
\begin{enumerate}
\item \textbf{Model One}: $\theta_i \overset{\text{i.i.d}}{\sim} Beta(0, 0)$
and $X_i \mid \theta_i \overset{\text{ind}}{\sim} Bin(n_i,
\theta_i)$. This model uses an uniformative prior on $\theta_i$, and
the posterior mean estimate of each $\theta_i$ coincides with the
frequentist MLE $X_i/n_i$. This model performs poorly for the present
dataset. In particular, rankings of counties by estimated risk (e.g.,
the top 100 values of $\hat{\theta}_i$) are 
dominated by counties with small populations. Moreover, the model
yields inaccurate predictions for future county-level death counts.
\item \textbf{Model Two}: $\theta_i \overset{\text{i.i.d}}{\sim} Beta(15, 150000)$
and $X_i \mid \theta_i \overset{\text{ind}}{\sim} Bin(n_i,
\theta_i)$. This model employs a highly informative prior on
$\theta_i$, where the hyperparameters $(15, 150000)$ may be viewed as
arising from a Beta distribution fitted to historical data. The prior
encodes strong prior knowledge about the prevalence of kidney cancer
mortality. This leads to substantial shrinkage of the posterior
estimates towards the prior mean. Unsurprisingly, given the strength
of the prior information, this model yields good empirical performance
on the dataset. 
\item \textbf{Model Three}: $\log a, \log b
  \overset{\text{i.i.d}}{\sim} \text{uniform}(-\infty, \infty)$,
  $\theta_i   \overset{\text{i.i.d}}{\sim} Beta(a, b)$, $X_i \mid
  \theta_i \overset{\text{ind}}{\sim} Bin(n_i, \theta_i)$. This model
  places an uninformative prior on the hyperparameters $a$ and $b$,
  and hence does not inject strong prior knowledge about the
  prevalence of kidney cancer mortality. Compared to Models One and
  Two, this is a substantially more flexible and principled model, as
  it allows the amount of shrinkage to be learned from the data rather
  than fixed a priori.

  When fitted to the data, the posterior distribution of $(a, b)$
  concentrates sharply around values close to $(15, 150000)$ (see the
  code file for this lecture), indicating that the model successfully
  infers from the data that kidney cancer deaths are rate events. As a
  consequence, the posterior estimates of the individual $\theta_i$
  are adaptively shrunk toward a common mean, with the degree of
  shrinkage depending on the corresponding population sizes
  $n_i$. This hierarchical borrowing of strength leads to stable
  estimation, sensible county-level rankings, and accurate predictions
  of future death counts, and overall the model performs very well.

  Note that this model does not require any a priori knowledge
  specific to kidney cancer and is therefore broadly applicable. For
  example, the same hierarchical framework can be used to predict
  batting averages in the baseball example, which will be explored in
  Homework Two. 
\end{enumerate}

\section{Normal Likelihoods}
We will next discuss Bayesian inference with normal likelihoods, and
also study some connections to the James-Stein estimator. Normal
likelihoods are applicable even in Binomial situations, such as in the
baseball data analysis example from \citet[Chapter 1]{Efron}. Here
$X_i$ denotes the number of hits in $n_i = n = 45$ at bats for player
$i$. The natural model is: $X_i \sim Bin(n_i, p_i)$ with parameter
$p_i$ but a normal likelihood (as opposed to Binomial likelihood) can
also be used because by invoking the Central Limit Theorem, we can
write
\begin{align*}
  \frac{X_i}{n_i} \overset{\text{approx}}{\sim} N\left(p_i,
  \frac{p_i(1-p_i)}{n_i} \right). 
\end{align*}
This is because:
\begin{align}\label{clt}
  \sqrt{n} \left(\frac{X_i}{n_i} - p_i \right)
  \overset{\text{Law}}{\rightarrow} N(0, p_i(1 - p_i)) \text{  as } n
  \rightarrow \infty. 
\end{align}
Note that the unknown parameter $p_i$ appears in the variance
above. If we want to work with normal likelihoods with known variance,
a clean way is to use a variance stabilizing transformation. This is
justified by the use of Delta method.
\begin{theorem}[Delta Method]
  If $\sqrt{n}(T_n - p) \overset{\text{Law}}{\rightarrow} N(0,
  \tau^2)$ as $n \rightarrow \infty$, then
  \begin{equation*}
    \sqrt{n} \left(g(T_n) - g(p) \right)
    \overset{\text{Law}}{\rightarrow} N(0, \tau^2 (g'(\theta))^2) 
  \end{equation*}
  as $n \rightarrow \infty$ provided $g'(p)$ exists and is
  non-zero.  
\end{theorem}
Informally, the Delta method states that if $T_n$ has a limiting
Normal distribution, then $g(T_n)$ also has  a limiting normal
distribution and also gives an explicit  formula for the asymptotic
variance of $g(T_n)$. This is surprising because $g$ can be linear or
non-linear. In general, non-linear functions of normal random
variables do not have a normal distribution. But the Delta method
works because under the assumption that $\sqrt{n}(T_n - p) \overset{\text{Law}}{\rightarrow}
N(0, \tau^2)$, it follows that $T_n$  converges in probability to
$p$ so that $T_n$ will be close to $p$ at least for large $n$. In a neighborhood of
$p$, the non-linear function $g$ can be approximated by a linear
function which means that $g$ effectively behaves like a linear
function. Indeed, the Delta method is a consequence of the
approximation: 
\begin{equation*}
  g(T_n) - g(p) \approx g'(p) \left(T_n - p \right). 
\end{equation*}
By the Delta method and \eqref{clt}, we have
\begin{equation*}
  \sqrt{n_i} \left(g(X_i/n_i) - g(p_i) \right)
  \overset{\text{Law}}{\rightarrow} N(0, (g'(p_i))^2 p_i(1 -
  p_i)). 
\end{equation*}
The variance above will not depend on $p_i$ if we choose the
function $g$ so that
\begin{equation*}
  g'(p_i) = \frac{1}{\sqrt{p_i(1 - p_i)}}
\end{equation*}
Solving this for $g$, we get $g(p_i) = 2
\arcsin(\sqrt{p_i})$. This is a variance stabilizing transformation
for the Binomial. Thus by the Delta method, we have
\begin{equation*}
  2 \sqrt{n} \left(\arcsin(\sqrt{X_i/n_i}) - \arcsin(\sqrt{p_i})
  \right)   \overset{\text{Law}}{\rightarrow} N(0, 1). 
\end{equation*}
While working with binomial counts, it thus makes sense to transform
the observation $X_i$ and parameter $p_i$ into
\begin{equation*}
  Y_i := 2 \sqrt{n_i} \arcsin(\sqrt{X_i/n_i}) ~~ \text{ and } ~~
  \theta_i := 2 \sqrt{n_i} \arcsin(\sqrt{p_i})  
\end{equation*}
respectively and then work with the likelihood
\begin{equation*}
  Y_i | \theta_i \sim N(\theta_i, 1). 
\end{equation*}
We will study Bayesian estimation under this model in the next
lecture. 


\begin{thebibliography}{99}

  
%\bibitem[Fisher(1934)]{Fisher1934}
%R.~A. Fisher,
%\newblock ``Probability, likelihood and quantity of information in the logic of
%uncertain inference,''
%\newblock \emph{Proceedings of the Royal Society of London. Series~A,
%Containing Papers of a Mathematical and Physical Character}, vol.~146,
%no.~856, pp.~1--8, 1934.

\bibitem[Efron(2010)]{Efron}
Efron, B. (2010)
\textit{Large-scale inference: empirical Bayes methods for estimation,
testing, and prediction}. Cambridge University Press, 2010.
  
%     \bibitem[Jaynes(2003)]{Jaynes} Jaynes, E. T, (2003)
%  \textit{Probability theory: the logic of science}. Cambridge
%  University Press, 2003.

%  \bibitem[Sinai(1992)]{Sinai} Sinai, Y. G, (1992)
%  \textit{Probability theory: an introductory course}. Springer
%  Verlag, 1992.  

%\bibitem[{deGroot(1986)}]{DeGrootBlackwell}DeGroot, M. H. (1986)
%       A conversation with David Blackwell.
%\newblock \emph{Statistical Science}, \textbf{1}, 40--53.

%         \bibitem[Mosteller(1987)]{Mosteller} Mosteller, F, (1987)
%  \textit{Fifty challenging problems in probability with
%    solutions}. Courier Corporation, 1987.

%    \bibitem[MacKay(2003)]{MacKay} MacKay, D. J. C., (2003)
%  \textit{Information theory, inference, and learning
%  algorithms}. Cambridge University Press, 2003.

%\bibitem[{Lindley(2013)}]{Lindley}Lindley, D. V. (2013)
%   \textit{Understanding Uncertainty}. John Wiley and sons, 2013.


\end{thebibliography}







\end{document}

%%% Local Variables:
%%% mode: latex
%%% TeX-master: t
%%% End:
